\subsection{Taxonomy}
\label{sec:mit-taxonomy}

We taxonomize mitigations across four layers, from closest to the
parser to closest to the wire. Table~\ref{tab:mitigations}
summarizes coverage per layer; the prose below treats each
layer's strengths and limits.

\begin{table}[h]
  \centering
  \footnotesize
  \caption{Mitigation taxonomy. \emph{Default-deployed} means the
  mitigation is on without any operator action; \emph{available}
  means an operator can enable it but does not by default;
  \emph{absent} means no equivalent primitive exists in the
  ecosystem.}
  \label{tab:mitigations}
  \begin{tabularx}{\linewidth}{@{}l X l l l@{}}
    \toprule
    Layer & Mechanism & Default-deployed & Available & Absent \\
    \midrule
    Server-side aggregation
      & coalesce N chunk events into one handler delivery
      & Kestrel (Pipelines)
      & \fwid{HttpObjectAggregator} (Netty)
      & all others \\
    Application-side chunk limit
      & per-request cap on chunk count
      & \emph{none}
      & \fname{BodyExt::limit\_chunks(N)} (hyper)
      & all others \\
    Application-side byte limit
      & total body byte cap (\fname{Content-Length} or otherwise)
      & most servers
      & N/A
      & N/A \\
    Frontend buffering
      & reverse proxy aggregates body before forwarding
      & nginx (\fname{proxy\_request\_buffering on}), Apache
        (\fname{mod\_proxy\_http})
      & N/A
      & HAProxy (streams by default) \\
    Transport-layer rate limit
      & per-IP packet/sec cap
      & N/A (operator opt-in)
      & \fname{iptables}, \fname{nftables}, BPF
      & N/A \\
    \bottomrule
  \end{tabularx}
\end{table}

\subsection{Server-side aggregation}
\label{sec:mit-server}

The structural fix is to coalesce consecutive chunk events into a
single application delivery, as Kestrel does (Section~\ref{sec:rc-kestrel}).
Netty exposes the
\fname{HttpObjectAggregator} pipeline handler precisely for this
purpose~\cite{netty_aggregator}: it accumulates
\fname{HttpContent} fragments up to a configured byte limit and
fires a single \fname{FullHttpRequest} event downstream. The
mechanism exists; what is missing in Vert.x, Spring Boot, and
every Netty-derivative we tested is that the aggregator is
\emph{not installed by default} because doing so would change the
streaming semantics that the framework's public API contracts
guarantee.

In other words: the cost of correctness in this class is breaking
the streaming contract by default. Frameworks that have made the
opposite choice (Pipelines in Kestrel; full buffering in Puma's
Tempfile handler boundary) closed the bug class at the price of
that contract. Frameworks that preserve streaming pay the
per-chunk dispatch cost for every chunk forever.

We argue (Section~\ref{sec:recommendations}) that the right
default for new frameworks is to aggregate up to a configurable
threshold (e.g., \SI{64}{\kibi\byte} or
\SI{100}{\milli\second} since last flush, whichever first), and
expose streaming as an explicit opt-in.

\subsection{Application-side chunk limit}
\label{sec:mit-app}

\fwid{hyper} closed issue \#4008~\cite{hyper_4008} as
\fname{not\_planned} on 2026-01-12 with the explicit position that
chunked-TE amplification is application responsibility, and pointed
users at \fname{http\_body\_util::Limited} (for byte caps) and the
proposed \fname{BodyExt::limit\_chunks(N)} helper (for chunk-count
caps; not yet shipped in
\fname{http\_body\_util}~\cite{hyper_4008}). No other server in
our sample exposes an equivalent chunk-count cap. Application code
in any other ecosystem that wants to defend against this attack
must write its own middleware that counts chunks during body
read.

This is functional but not pleasant: every application has to
reinvent the cap, and many will not. The hyper maintainer position
is internally consistent (server-side aggregation breaks
streaming; chunk caps are an app-level concern) but does not scale
across an ecosystem where most application code is unaware of the
issue.

\subsection{Frontend buffering}
\label{sec:mit-frontend}

The most widely-given advice ("deploy behind nginx") turns out to
be valid for nginx and Apache, not for HAProxy.

\paragraph{nginx.}
\fname{proxy\_request\_buffering on} is the default. With it
enabled, nginx reads the entire decoded request body into
\fname{client\_body\_buffer\_size} (default 8--16 KB; spills to
disk via \fname{client\_body\_temp\_path} beyond that). Only
after the body is fully received does nginx open the upstream
connection and forward the body in a single \fname{Content-Length}
request. The upstream sees one well-formed request, not a chunked
stream. Per-chunk decoder cost is paid by nginx, not by the
upstream.

We verified this empirically: an instrumented Python upstream
behind \fwid{nginx} 1.27-alpine (default \fname{proxy\_request\_buffering}
\fname{on}) received an attacker chunked-TE body of $N{=}\num{50000}$
one-byte chunks as a single 50\,165-byte \fname{recv()} carrying a
\fname{Content-Length: 50000} header and no
\fname{Transfer-Encoding} header. The same probe against the same
upstream with \fname{proxy\_request\_buffering~off} was forwarded
as the original chunked stream (five separate \fname{recv()}s, no
\fname{Content-Length}, \fname{Transfer-Encoding: chunked}
preserved). The full reproducer is at
\texttt{wave2/c/nginx-as-proxy/} and the captured upstream logs
at \texttt{wave2/c/nginx-as-proxy/logs/scenario-\{A,B\}-*.log}.

For every Python, Node, Ruby, Java, or .NET application deployed
behind a default nginx, this is full mitigation of the per-chunk
amplification class for the upstream. Nginx itself still pays the
parsing cost (Section~\ref{sec:synth-matrix} reports a synthetic
nginx-as-origin cost of $\sim$\SI{114}{\micro\second}/chunk under
an OpenResty Lua sink, among the highest per-chunk cells we
measured) but nginx is a hardened C server with well-tuned
worker concurrency, and as a buffering proxy it never delivers the
body per-chunk to the upstream at all, so the impact is bounded.

\paragraph{Apache httpd.}
\fname{mod\_proxy\_http} behaves the same way:
\fname{ProxyIOBufferSize} (default 8192) bounds per-request
buffering. Upstream sees a coalesced body.

\paragraph{HAProxy.}
HAProxy streams by default. \fname{http-buffer-request} can be
enabled but only waits for one buffer's worth of bytes
(\fname{tune.bufsize}, default \SI{16}{\kibi\byte}), not a full
per-chunk accumulator. Upstream servers behind HAProxy receive
chunked-TE requests directly and inherit the per-chunk cost. We
have not seen this clearly stated anywhere in the HAProxy
documentation; operators relying on "HAProxy is in front of my
app" as a mitigation are mistaken.

\paragraph{Cloud frontends.}
We have not measured Cloudflare, Fastly, AWS ALB, GCP HTTP(S) LB,
or Akamai. Each of these CDNs and load balancers has its own
defaults for chunked-body buffering. We flag this as future work
(Section~\ref{sec:future}). Operators should empirically verify
whether their frontend buffers chunks before assuming a mitigation
is in place.

\subsection{Transport-layer rate limiting}
\label{sec:mit-transport}

A coarser fallback: rate-limit incoming TCP packets per source IP
or per connection. This is universally available (\fname{iptables},
\fname{nftables}, eBPF, dedicated DDoS appliances) but expensive
to operate correctly:

\begin{itemize}
  \item Slow legitimate uploads (cellular networks, satellite,
    constrained IoT) will be caught by aggressive packet-rate
    caps.
  \item Per-IP rate limiting fails open under distributed attacks
    (botnets, residential proxies).
  \item Per-connection caps catch low-bandwidth slow-drip but
    not high-bandwidth bursts.
\end{itemize}

We recommend transport-layer rate limiting only as a defense
in depth, not as the primary mitigation.

\subsection{Summary recommendation per deployment tier}
\label{sec:mit-summary}

\begin{description}
  \item[Direct-exposed application server.] If you cannot deploy
    a frontend, you must apply Server-side aggregation
    (Kestrel only) or Application-side chunk limits (hyper, or
    custom middleware everywhere else).
  \item[Behind nginx or Apache (default config).] Mitigated for the
    upstream. Verify
    \fname{proxy\_request\_buffering on} (nginx) or
    \fname{ProxyIOBufferSize} default (Apache).
  \item[Behind HAProxy.] Not mitigated; you must apply
    application-side caps or enable a body-buffering proxy in
    front of HAProxy.
  \item[Behind a CDN.] Verify with the CDN provider that
    chunked-TE bodies are aggregated before being forwarded to
    your origin. Most providers do not document this behavior
    explicitly.
\end{description}
